\section{Conclusions and perspectives}\label{sec:conclusion_and_future_direction}

We have studied the subfield of machine learning with rejection and provided a higher-level overview of existing research. To conclude, we revisit our key research questions and point to new directions that future research might take.

\subsection{Research questions revisited}
This survey paper is built around eight key research questions, introduced in the introduction. In this section, we revisit each of these questions and briefly summarize our findings.

\paragraph{How can we formalize the conditions for which a model should abstain from making a prediction?}
In Section~\ref{sec:preliminaries}, we identify two types of rejections: ambiguity rejection and novelty rejection.
Ambiguity rejection abstains from making a prediction an example falls in a region where the target value is ambiguous (e.g., close to the decision boundary in classification tasks).
This could be due to a non-deterministic true relation between the features and target variable, or due to a hypothesis space that is not able to capture the true relation. 
Novelty rejection abstains from making a prediction on examples that are rare with respect to the given training set. For such an example, there is no guarantee that the model correctly extrapolates to this untrained region, making it likely that the model mispredicts the example.

\paragraph{How can we evaluate the performance of a model with rejection?}
Standard machine learning evaluation is focused on a model's predictive quality. However, in machine learning with rejection, there exists a trade-off between the predictive quality and the proportion of rejected examples.

In Section~\ref{sec:metrics}, we provide an overview of techniques evaluating both the prediction and rejection quality of models with rejection. We identify three categories: metrics evaluating models with a given rejection rate, metrics evaluating the overall model performance/rejection trade-off, and metrics evaluating models through a cost function.

\paragraph{What architectures are possible for operationalizing (i.e., putting this into practice) the ability to abstain from making a prediction?}
We categorize machine learning methodologies with rejection in three different architectures, depending on the relationship between the predictor and the rejector: separated, dependent and integrated rejector. These categories are introduced and mapped to the existing literature in Sections~\ref{sec:separated_rejector},~\ref{sec:dependent_rejector} and~\ref{sec:integrated_rejector}. 

\paragraph{How do we learn models with rejection?}
For each architecture, we discuss the main techniques to learn a model with rejection and related these to the existing literature in Sections~\ref{sec:separated_rejector},~\ref{sec:dependent_rejector} and~\ref{sec:integrated_rejector}. 
First, the separated rejector is usually learned independently of the predictor.
Second, learning the dependent rejector entails learning for which examples the predictor is likely to mispredict using a confidence function. Both architectures need setting a rejection threshold.
Third, integrated rejector needs a unique algorithm for learning predictor and rejector in tandem. Usually, this architecture relies on designing an objective function.

\paragraph{What are the main pros and cons of using a specific architecture?}
Each architecture, discussed in Sections~\ref{sec:separated_rejector},\ref{sec:dependent_rejector}, and~\ref{sec:integrated_rejector}, offers distinct benefits and drawbacks. Separated rejectors show broad applicability, as they can be combined with any predictor. However, they often yield sub-optimal rejection performance since they do not learn from the predictor's mispredictions. On the other hand, dependent rejectors have reduced, yet still high, applicability, relying on a specific confidence function learned from the predictor's output, but they can enhance the rejection quality by leveraging the predictor's mispredictions. Finally, integrated rejectors necessitate joint design with the predictor, but learning a single model for prediction and rejection improves the overall performance for both prediction and rejection tasks.

\paragraph{How can we combine multiple rejectors?}
We discuss the combination of multiple rejectors for enabling various types of rejections within a unique model. There are two approaches for combining rejectors. First, when rejectors do not overlap, a logical ``or'' rule is applied, rejecting an example if any of the rejectors rejects it. Second, when rejectors overlap in some regions and disagree on the type of rejection, a multi-step architecture is used, ordering rejectors by importance to make decisions based on the most relevant rejector.

\paragraph{Where does the need for machine learning with rejection methods arise in real-world applications?}
On a high level, machine learning with rejection is typically used in applications where incorrect decisions can have severe consequences, both financially and safety-related. These consequences motivate the need for robust and trustworthy machine learning. In Section~\ref{sec:applications}, we provide an overview of application areas in which machine learning with rejection is already used.

\paragraph{How does machine learning with rejection relate to other research areas?}
Section~\ref{sec:related_research_areas} shows that machine learning with rejection is closely related to several other subfields of machine learning. This relation sometimes leads to terminology and techniques overlapping or inspired by these other domains. In contrast, other cases show machine learning with rejection from a broader perspective. In this survey, we related machine learning with rejection to uncertainty quantification, anomaly detection, active learning, class-incremental learning, delegating classifiers, and meta-learning.

\subsection{Future directions}
Given its significance for the usage of machine learning in real-world problems and the growing attention for trustworthy AI, we expect machine learning with rejection to remain an active research field. In this section, we briefly discuss three key research directions for which we see a strong need.

\paragraph{Standard settings to compare different models with rejection.}
A large number of machine learning models with rejection already exist. However, these are typically evaluated on custom or even proprietary data. This makes it difficult to benchmark and compare the different approaches. While some papers use publicly available datasets, there is no standard benchmark set for machine learning with rejection.
Additionally, applying multiple strategies to evaluate the rejector offers a better view of an algorithm's performance and improves comparability.

\paragraph{Partial rejection for machine learning models.}
A promising avenue for further exploration is the concept of partial abstention. 
\changed{Nowadays, machine learning problems often involve seeking elaborated predictions rather than simple scalar or class values as in classification and regression tasks. For instance, in multi-label classification, a prediction for an instance is a subset of possible class labels.
In such cases, the idea of abstaining from a complete prediction can be extended to partial abstention, where the learner delivers predictions on \emph{some but not necessarily all}
class labels, according to its level of certainty~\citep{Nguyen2020}. 
}
This has the key benefit of providing a middle ground between making predictions for the entire structure and completely abstaining from making any predictions.


\paragraph{Algorithms enabling models with rejection in domains other than classification.}
Most papers on machine learning with rejection study supervised classification problems. Modern machine learning tackles numerous other tasks such as regression, forecasting, and clustering, or even semi-supervised and self-supervised feedback loops. We believe that the rejection can also be of use in such areas. However, since only a handful of relevant studies exist, this requires more attention from the research community. 
Future research can also focus on integrating rejection-related variables into statistical frameworks utilized in educational measurement, such as Item Response Theory (IRT). This integration has the potential to improve the precision of item calibration, trait estimation, and the interpretation of test scores. Furthermore, it can provide valuable insights into the psychological aspects of learning, enabling the development of more precise instructional strategies and interventions.
